\documentclass[../dissertacao]{subfiles}

\begin{document}
\chapter{Primeiro capítulo}

\section{Aqui está uma seção}

\begin{definition}[Espaços Completos]\label{def:completo}
    Um espaço métrico é \emph{completo} se toda sequência de Cauchy converge para um elemento do espaço. 
\end{definition}

Podemos referenciar a \zcref{def:completo} desta forma. Note que o comando já adiciona a palavra "Definição". Caso não queira esse comportamento, use \zcref[noname]{def:completo}
Aqui vai mais alguns exemplos do 'zcref':

\begin{enumerate}
    \item primeiro item \label{item:1}
    \item segundo item \label{item:2}
    \item terceiro item \label{item:3}
\end{enumerate}

Confira os \zcref{item:1,item:2,item:3}.

\begin{theorem}\label{thm:a}
    teorema...
\end{theorem}
\begin{theorem}\label{thm:b}
    teorema...
\end{theorem}
\begin{theorem}\label{thm:c}
    teorema...
\end{theorem}

\zcref{thm:a,thm:c} e também o \zcref{thm:b}

\lipsum

\end{document}